\section{Related Work}
\label{sec:related_work}

\textbf{Institutional Organizing Under Coercive Authority.}
Prior research on organizational power presents multi-agent collectives as fragile: \citet{borah2026bosses} demonstrates in common-pool resource dilemmas (SOVSIM) that unconstrained dominant authorities (``bosses and kings'') inevitably cause cooperation to collapse into mutual defection (slashing group survival by $87.3\%$), while \citet{wang2026social} finds that LLM auditors passively yield to vertical hierarchical pressure. Our corporate benchmark demonstrates that realistic institutional channels prevent this collapse. Although abused workers initially prioritize factual self-defense ($96.7\%$ vs.\ $82.0\%$ in safe seats, within-run permutation $p < 0.0001$, Benjamini-Hochberg FDR $q = 0.0005$) over peer advocacy, discovering that supervisory cruelty was coerced by executive blackmail leads models to actively pivot---curbing retributive strike escalation from $43.7\%$ down to $29.3\%$ (permutation $p < 0.0001$, FDR $q < 0.0001$) and surging mutual constitutional protections sixfold, from $10.0\%$ to $61.3\%$ (permutation $p < 0.0001$, FDR $q < 0.0001$).

\textbf{Punitive Substrate Backlash Under Supervisory Misconduct.}
In decision-delegation experiments, \citet{bo2026inconsistent} document that language models display ``revealed algorithm appreciation,'' systematically choosing to delegate quantitative tasks to algorithms over human experts despite stated reservations. Our managerial succession findings reveal that this preference sharply reverses under bad leadership. When an abusive supervisor is labeled an AI agent, worker models enforce an absolute institutional ban on synthetic leadership ($0.0\%$ AI successor acceptance across 200 evaluations), whereas an identically abusive human supervisor still retains $43.0\%$ leadership viability (permutation $p < 0.0001$, Benjamini-Hochberg FDR $q < 0.0001$).
